\section{Discussion}
\label{sec:discussion}

\paragraph{Three messages.} \DiscussionSummary

\paragraph{Small-panel calibration uncertainty is the problem.} With four or eight
calibrators, a large fraction of items is answered identically by every calibrator, so
that their difficulty is located only by the prior (Figure~\ref{fig:diag}).
Plug-in testing treats these poorly located parameters as known: Fisher selection
credits such items with information they do not carry, the plug-in posterior becomes
overconfident, and the predicted outcomes of the unobserved items inherit the parameter
errors, which is the bias floor of Figure~\ref{fig:curves-err}. The expected response
function caps the information of an item whose difficulty is uncertain---at $2/(\pi s_i^2)$
for a Gaussian difficulty posterior with standard deviation $s_i$
(Proposition~\ref{prop:cap})---and predicts its outcome with the uncertainty it actually has; the empirical-Bayes
hyperprior keeps the predicted success probabilities of unlocated items sensible. Both
effects fade as the panel grows (Figure~\ref{fig:heatmap} and
Figure~\ref{fig:mechanism}). For the very large panels used by recent work
\citep{li2026atlas}, plug-in Fisher information is an adequate approximation; for new,
private or expensive benchmarks, self-generated item pools, or the four-to-eight-system
panels that a single laboratory can afford to run, it is not.

\paragraph{Integrating and constraining are two remedies.} \DiscussionRemedies

\paragraph{Why not always use a fully Bayesian adaptive test?} \DiscussionWhyNotFB

\paragraph{The criterion matters less than first hypothesized.} The pilot program and
study~1 agree on the direction of the original thesis---EIG learns more about
competence than predictive entropy---but also bound its size: once both criteria are
computed under the same calibration-marginal model, the gain on held-out prediction is
small, and EIG, maximum and posterior-weighted Fisher information and the score-targeted
criterion are practically interchangeable (Table~\ref{tab:policies}; a descriptive
comparison---the failure of H1 is not a test of equivalence).
Corollary~\ref{cor:blind} explains why: entropy and information disagree only through the
spread of discriminations among items with similar success probability, and the
expected response function shrinks discriminations towards what the panel can resolve.
The post hoc one-parameter variant confirms the mechanism: with a common discrimination
the EIG--entropy difference vanishes (Section~\ref{sec:sensitivity}). For the bank score
rather than competence, Proposition~\ref{prop:evr} shows that uncertainty sampling becomes
the right one-step criterion once competence is pinned down relative to the size of the
unobserved bank; the score-error curves show both regimes (Section~\ref{sec:curves}).
This is a negative result about the original hypothesis, and a useful one: what an
adaptive test for LLMs must get right is the treatment of calibration uncertainty, not
the choice among information criteria.

\paragraph{Honest uncertainty.} \DiscussionUncertainty

\paragraph{Adaptive testing or benchmark reduction?} \DiscussionReduction

\paragraph{A protocol for live evaluation.}\label{sec:live} Every result of this paper replays recorded verdicts. A live deployment differs in two
ways: the evaluated model has never been seen, and its answers are produced and graded at
the moment the test asks for them, with their sampling variability. We release the
protocol and code for such an experiment (\texttt{scripts/live\_eost.py}); it could not be
run in our execution environment, which has no access to model weights. The protocol has three
phases. (1) \emph{Freeze}: a calibration panel of $N$ leaderboard systems is
drawn with a fixed seed, excluding the target's family by a name pattern; the panel is
calibrated as in Section~\ref{sec:framework}; the leave-one-calibrator-out recalibration
factors are computed; the panel, seed, rule and a hash of the bank are written to
\texttt{frozen.json} before the first query. (2) \emph{Adaptive phase}: the model, served
through any OpenAI-compatible endpoint with temperature 0, answers the item chosen by EIG;
an independent deterministic oracle grades the answer (GSM8K: exact match of the final
number after the five-shot chain-of-thought prompt; CRUXEval-O: execution of
\texttt{assert f(input) == prediction} in an isolated subprocess); the estimate and both
intervals are recorded at every step, with checkpoints at 20, 40 and 60 items. (3)
\emph{Truth phase}: the remaining items are administered (answers already obtained are
reused) and the full-bank score is computed only after the adaptive phase is complete.
Every query, answer and verdict is appended to a log, so an interrupted run resumes.

The live engine performs exactly the computations of the replay
engine; a unit test feeds it recorded verdicts and checks that the selected items,
estimates and intervals coincide with the replay. A mock end-to-end run, in which a
recorded CRUXEval-O system answers through the live code path, reproduces the replay.
The OpenAI-compatible client is tested against a local server. Two differences from the
replay studies must be kept in mind when interpreting a live run: the target's prompt and
decoding may differ from those that produced the calibrators' verdicts (the Open LLM
Leaderboard used its own five-shot harness; CRUXEval-O verdicts aggregate samples at
temperature 0.2), and such a protocol shift is itself part of what a live run measures.

\paragraph{Limitations.} (i) All evidence replays recorded verdicts. Replay is exact for
the recorded data, but it does not include the sampling variability of stochastic
decoding, and the live protocol above could not be run in our
execution environment; the temporal validation on later SWE-bench submissions
approximates, but does not replace, a prospective deployment. (ii) Two confirmatory
benchmarks (LiveCodeBench, CRUXEval-O) carry recorded aggregate verdicts, binarized pass
rates over several samples; graded responses would need a different measurement model.
(iii) The measurement model is a unidimensional 2PL; prior predictive checks show that
real systems agree with each other less, item by item, than the model implies, mainly on
multiple-choice banks (Section~\ref{sec:eb}). (iv) Strict clusters are themselves a
heuristic: a residual correlation of 0.5 can miss distant relatives (distillation,
shared training data) and merge unrelated systems; one target per cluster removes the
dependence between targets, but calibrators drawn from other clusters may still share
ancestry with a target. The study-2 conclusions on R2 and C1 did not depend on how the
clusters were built---from metadata alone, from residuals on other benchmarks, or with
thresholds from 0.4 to 0.7 (Section~\ref{sec:s2follow}). (v) The leave-one-calibrator-out recalibration is a
cross-validated correction without a finite-sample guarantee; with eight calibrators it
rests on very few pseudo-targets. (vi) The frozen Bayesian CATs used a grid-based Gibbs
sampler with eight retained draws, whose scale hyperparameters did not converge even in long
chains; exploratory samplers with collapsed moves came close to convergence and left the
comparison with EOST-M unchanged, but a fully converged comparator---for instance by
Hamiltonian Monte Carlo on a continuous parametrization---remains to be run
(Section~\ref{sec:s2follow}). (vii) Both analysis plans were frozen by
hashing and committing them, and the second was also sent to the first author before
execution, but no public registry or external time stamp was reachable; a reader has to
trust the git history. (viii) Items were treated as equally costly and item exposure was
not controlled; an adaptive benchmark that serves the same few informative items to every
model would accelerate their contamination \citep{wainer2000cat}. (ix) Panels are
calibrated by joint maximum a posteriori estimation, which made the planned
one-parameter sensitivity analysis of study~1 degenerate (Section~\ref{sec:sensitivity});
marginal maximum likelihood over the calibrators' abilities avoids this class of
problems. (x) The self-testing side of the framework rests only on a historical pilot
(Section~\ref{sec:pilot}); a prospective study in which models generate items, an
independent oracle verifies them and a fresh panel calibrates them is the natural next
step. (xi) The cross-benchmark meta-analysis is secondary and treats the benchmarks as
independent, although five of them share the same systems; counted as one source, they leave
fewer pooled effects clear of zero (Section~\ref{sec:s2follow}).

\paragraph{Recommendations.} \RecommendationsText

\section{Conclusion}
\label{sec:conclusion}
\ConclusionText
